% !TeX program = pdflatex
% !TeX encoding = UTF-8
\documentclass[12pt,aspectratio=169]{beamer}
% SimulationStyle.tex
% Standard academic Beamer style for the Science Quiz Simulation.
% Based on the conventional Madrid Beamer layout used in the Seoul High School lecture deck.

\usetheme{Madrid}
\useinnertheme{rounded}
\usefonttheme{professionalfonts}

\usepackage{kotex}
\usepackage[scaled=0.94]{helvet}
\renewcommand{\familydefault}{\sfdefault}
\usepackage{amsmath,amssymb,mathtools}
\usepackage[version=4]{mhchem}
\usepackage{xcolor}

\definecolor{POSTECHRed}{RGB}{200,16,46}
\definecolor{POSTECHDark}{RGB}{25,25,25}

\setbeamercolor{normal text}{fg=black,bg=white}
\setbeamercolor{structure}{fg=POSTECHRed}
\setbeamercolor{alerted text}{fg=POSTECHRed}
\setbeamercolor{palette primary}{fg=white,bg=POSTECHRed}
\setbeamercolor{palette secondary}{fg=white,bg=POSTECHDark}
\setbeamercolor{palette tertiary}{fg=white,bg=POSTECHRed}
\setbeamercolor{palette quaternary}{fg=white,bg=POSTECHDark}
\setbeamercolor{frametitle}{fg=white,bg=POSTECHRed}
\setbeamercolor{title}{fg=white,bg=POSTECHRed}
\setbeamercolor{titlelike}{fg=white,bg=POSTECHRed}
\setbeamercolor{block title}{fg=white,bg=POSTECHRed}
\setbeamercolor{block body}{fg=black,bg=black!5}
\setbeamercolor{item projected}{fg=white,bg=POSTECHRed}
\setbeamercolor{subitem projected}{fg=white,bg=POSTECHRed}

\setbeamertemplate{navigation symbols}{}
\setbeamertemplate{items}[circle]
\setbeamertemplate{blocks}[rounded][shadow=false]
\setlength{\parskip}{0.22em}
\setbeamerfont{itemize/enumerate body}{size=\normalsize}
\setbeamerfont{itemize/enumerate subbody}{size=\small}
\setbeamerfont{block body}{size=\normalsize}

% Minimal three-part footline, following the same academic Beamer convention.
\makeatletter
\setbeamertemplate{footline}{%
  \leavevmode%
  \hbox{%
    \begin{beamercolorbox}[wd=.36\paperwidth,ht=2.25ex,dp=1ex,leftskip=1.2em]{author in head/foot}%
      \usebeamerfont{author in head/foot}20260917 Simulation
    \end{beamercolorbox}%
    \begin{beamercolorbox}[wd=.28\paperwidth,ht=2.25ex,dp=1ex,center]{title in head/foot}%
      \usebeamerfont{title in head/foot}POSTECH
    \end{beamercolorbox}%
    \begin{beamercolorbox}[wd=.36\paperwidth,ht=2.25ex,dp=1ex,rightskip=1.2em plus1fil]{date in head/foot}%
      \usebeamerfont{date in head/foot}\insertframenumber{} / \inserttotalframenumber
    \end{beamercolorbox}%
  }%
  \vskip0pt%
}
\makeatother

\newenvironment{problemframe}[2]{%
  \begin{frame}[t]{Problem #1 -- #2}%
  \large
}{%
  \end{frame}%
}

\newenvironment{solutionframe}[2]{%
  \begin{frame}[t]{Solution #1 -- #2}%
  \normalsize
}{%
  \end{frame}%
}

\newcommand{\answerbox}[1]{%
  \begin{block}{정답}
    \centering
    {\large\bfseries #1}
  \end{block}
  \vspace{0.2em}%
}

\newcommand{\solutionlabel}{%
  \textbf{해설.}\hspace{0.35em}%
}

\newcommand{\sourceinfo}[2]{%
  \par\vspace{0.35em}%
  {\scriptsize\color{gray}출처: \href{#2}{#1}}%
}

\newenvironment{simchoices}{%
  \begin{enumerate}
    \setlength{\itemsep}{0.42em}%
    \setlength{\parskip}{0pt}%
}{%
  \end{enumerate}%
}

\newcommand{\dd}{\mathop{}\!\mathrm{d}}
\newcommand{\Q}{\mathbb{Q}}


\hypersetup{
  pdftitle={20260917 Science Quiz Simulation},
  pdfsubject={2026 KAIST-POSTECH Science War Science Quiz Simulation},
  pdfcreator={pdfLaTeX}
}

\title[Science Quiz Simulation]{Science Quiz Simulation}
\subtitle{2026 KAIST-POSTECH Science War}
\author{}
\institute{}
\date{2026.09.17}

\begin{document}

\begin{frame}[plain]
  \titlepage
\end{frame}

% 01 Physics - 문제
\begin{problemframe}{01}{Physics}
길이 $L$, 질량 $m$인 균일한 막대의 한쪽 끝을 마찰 없는 축으로 고정하였다. 막대를 수평으로 놓은 뒤 정지 상태에서 놓았을 때, 막대가 처음으로 수직 아래 방향을 지나는 순간의 각속력 $\omega$를 구하시오. 중력가속도는 $g$이고 공기저항은 무시한다.
\end{problemframe}

% 01 Physics - 정답
\begin{solutionframe}{01}{Physics}
\answerbox{$\displaystyle \omega=\sqrt{\frac{3g}{L}}$}
\solutionlabel
질량중심은 $L/2$만큼 내려가므로 잃은 위치에너지는 $mgL/2$이다. 한쪽 끝을 지나는 축에 대한 균일한 막대의 관성모멘트는 $I=mL^2/3$이므로
\[
\frac{mgL}{2}=\frac12 I\omega^2
=\frac16 mL^2\omega^2.
\]
따라서 $\omega^2=3g/L$이다. 핵심은 병진운동이 아니라 \textbf{고정축 회전의 에너지 보존}으로 보는 것이다.
\end{solutionframe}

% 02 Chemistry - 문제
\begin{problemframe}{02}{Chemistry}
$25\,^{\circ}\mathrm{C}$에서 다음 아연 농도전지가 있다.
\[
\ce{Zn(s) | Zn^{2+}(0.010 M) || Zn^{2+}(1.0 M) | Zn(s)}.
\]
활동도는 농도로 근사하고, $25\,^{\circ}\mathrm{C}$에서 $2.303RT/F=0.0592\,\mathrm{V}$라 하자. 이 전지의 전위차 $E_{\mathrm{cell}}$의 크기를 구하시오.
\end{problemframe}

% 02 Chemistry - 정답
\begin{solutionframe}{02}{Chemistry}
\answerbox{$\displaystyle E_{\mathrm{cell}}=0.0592\,\mathrm{V}$}
\solutionlabel
아연의 반쪽 반응은 $\ce{Zn^{2+}+2e^- -> Zn}$이므로 $n=2$이다. 농도전지에서는 높은 $\ce{Zn^{2+}}$ 농도 쪽이 환원전극이 되고
\[
E_{\mathrm{cell}}
=\frac{0.0592}{2}\log_{10}\!\left(\frac{1.0}{0.010}\right)
=\frac{0.0592}{2}\times2
=0.0592\,\mathrm{V}.
\]
같은 redox couple을 쓰므로 $E^\circ$가 상쇄된다는 점이 핵심이다.
\end{solutionframe}

% 03 Biology - 문제
\begin{problemframe}{03}{Biology}
100\%의 증폭 효율을 갖는 qPCR에서 같은 primer pair를 사용하였다. Sample A의 $C_t$가 18, Sample B의 $C_t$가 21이었다. 다른 조건이 모두 같을 때, 초기 target DNA 양 $N_A/N_B$를 구하시오.
\end{problemframe}

% 03 Biology - 정답
\begin{solutionframe}{03}{Biology}
\answerbox{$\displaystyle \frac{N_A}{N_B}=8$}
\solutionlabel
100\% 효율이면 한 cycle마다 target DNA가 두 배가 된다. 같은 threshold에 도달하는 데 A가 B보다 3 cycles 적게 필요했으므로
\[
\frac{N_A}{N_B}=2^{21-18}=2^3=8.
\]
즉 $C_t$가 작을수록 초기 template 양이 많다.
\end{solutionframe}

% 04 Mathematics - 문제
\begin{problemframe}{04}{Mathematics}
다음 $5\times5$ 행렬의 determinant를 구하시오.
\[
A=
\begin{pmatrix}
0&1&1&1&1\\
1&0&1&1&1\\
1&1&0&1&1\\
1&1&1&0&1\\
1&1&1&1&0
\end{pmatrix}.
\]
\end{problemframe}

% 04 Mathematics - 정답
\begin{solutionframe}{04}{Mathematics}
\answerbox{$\displaystyle \det A=4$}
\solutionlabel
모든 성분이 $1$인 행렬을 $J$라 하면 $A=J-I$. $J$의 eigenvalues는 $5$ 하나와 $0$ 네 개이므로 $A$의 eigenvalues는
\[
4,-1,-1,-1,-1.
\]
따라서
\[
\det A=4(-1)^4=4.
\]
직접 전개하기보다 \textbf{$J-I$ 구조와 eigenvalues}를 보는 것이 빠르다.
\end{solutionframe}

% 05 Computer Engineering - 문제
\begin{problemframe}{05}{Computer Engineering}
삭제 연산을 사용하지 않는 표준 Bloom filter에 대한 설명 중 옳은 것을 모두 고르시오.
\begin{simchoices}
\item false positive가 가능하다.
\item false negative가 가능하다.
\item ``없다''는 판정은 원소가 삽입되지 않았음을 보장한다.
\item ``있다''는 판정은 원소가 삽입되었음을 보장한다.
\end{simchoices}
\end{problemframe}

% 05 Computer Engineering - 정답
\begin{solutionframe}{05}{Computer Engineering}
\answerbox{(1), (3)}
\solutionlabel
Bloom filter는 여러 hash가 가리키는 bit를 $1$로 세우며, 표준형에서는 삽입된 원소의 bit가 다시 $0$으로 돌아가지 않는다. 따라서 \textbf{false positive는 가능하지만 false negative는 없다}. 그러므로 negative query는 확정적이지만 positive query는 확률적이다.
\end{solutionframe}

% 06 Physics - 문제
\begin{problemframe}{06}{Physics}
2025 Nobel Prize in Physics는 전기회로에서의 \textit{macroscopic quantum mechanical tunnelling}과 energy quantisation의 발견에 수여되었다. 해당 실험에서 두 초전도체 사이에 매우 얇은 절연층을 둔 핵심 소자의 이름을 쓰시오.
\end{problemframe}

% 06 Physics - 정답
\begin{solutionframe}{06}{Physics}
\answerbox{Josephson junction}
\solutionlabel
두 초전도체를 얇은 절연 장벽으로 분리하면 Cooper pair가 장벽을 quantum tunnelling할 수 있다. 이러한 구조가 \textbf{Josephson junction}이며, 2025 Nobel Physics의 macroscopic quantum circuit 실험의 핵심 요소였다.
\sourceinfo{Nobel Prize in Physics 2025}{https://www.nobelprize.org/prizes/physics/2025/press-release/}
\end{solutionframe}

% 07 Chemistry - 문제
\begin{problemframe}{07}{Chemistry}
다음 species 중 ground state에서 aromatic한 것을 모두 고르시오.
\begin{simchoices}
\item cyclopropenyl cation, $\ce{C3H3+}$
\item cyclobutadiene, $\ce{C4H4}$
\item cyclopentadienyl anion, $\ce{C5H5-}$
\item cyclooctatetraene, $\ce{C8H8}$
\end{simchoices}
\end{problemframe}

% 07 Chemistry - 정답
\begin{solutionframe}{07}{Chemistry}
\answerbox{(1), (3)}
\solutionlabel
Hückel rule에 따라 planar, cyclic, fully conjugated system이 $(4n+2)\pi$ electrons를 가지면 aromatic하다. $\ce{C3H3+}$는 $2\pi$ electrons, $\ce{C5H5-}$는 $6\pi$ electrons이므로 aromatic이다. Cyclobutadiene은 $4\pi$로 antiaromatic하고, cyclooctatetraene은 nonplanar tub conformation을 취해 nonaromatic하다.
\end{solutionframe}

% 08 Biology - 문제
\begin{problemframe}{08}{Biology}
2025 Nobel Prize in Physiology or Medicine는 peripheral immune tolerance에 관한 발견에 수여되었다. 사람에서 \textit{FOXP3}의 loss-of-function mutation이 가장 직접적으로 발달과 기능을 손상시키는 면역세포 집단은 무엇인가?
\end{problemframe}

% 08 Biology - 정답
\begin{solutionframe}{08}{Biology}
\answerbox{Regulatory T cell (Treg)}
\solutionlabel
FOXP3는 regulatory T cell의 발달과 억제 기능을 지휘하는 핵심 transcription factor이다. 기능성 FOXP3가 없으면 Treg-mediated peripheral tolerance가 무너져 심한 autoimmunity가 나타날 수 있다.
\sourceinfo{Nobel Prize in Physiology or Medicine 2025}{https://www.nobelprize.org/prizes/medicine/2025/press-release/}
\end{solutionframe}

% 09 Mathematics - 문제
\begin{problemframe}{09}{Mathematics}
다음 정적분의 값을 구하시오.
\[
I=\int_0^1\frac{x^3}{x^3+(1-x)^3}\,\dd x.
\]
\end{problemframe}

% 09 Mathematics - 정답
\begin{solutionframe}{09}{Mathematics}
\answerbox{$\displaystyle I=\frac12$}
\solutionlabel
$\displaystyle f(x)=\frac{x^3}{x^3+(1-x)^3}$라 두면 $f(x)+f(1-x)=1$이다. 따라서
\[
2I=\int_0^1\bigl(f(x)+f(1-x)\bigr)\,\dd x
=\int_0^1 1\,\dd x=1.
\]
그러므로 $I=1/2$. 적분을 전개하기 전에 \textbf{$x\leftrightarrow1-x$ symmetry}를 보는 것이 핵심이다.
\end{solutionframe}

% 10 Computer Engineering - 문제
\begin{problemframe}{10}{Computer Engineering}
단일 단계 page table을 사용하는 시스템에서 TLB lookup은 $10\,\mathrm{ns}$, main memory access는 $100\,\mathrm{ns}$이고 TLB hit rate는 $90\%$이다. Page fault는 없으며, TLB miss 시 page table entry를 memory에서 한 번 읽은 뒤 실제 data를 한 번 더 읽는다. Effective memory access time을 구하시오.
\end{problemframe}

% 10 Computer Engineering - 정답
\begin{solutionframe}{10}{Computer Engineering}
\answerbox{$\displaystyle 120\,\mathrm{ns}$}
\solutionlabel
TLB hit이면 $10+100=110\,\mathrm{ns}$, miss이면 $10+100+100=210\,\mathrm{ns}$가 걸린다. 따라서
\[
0.9(110)+0.1(210)=99+21=120\,\mathrm{ns}.
\]
Virtual memory 계산에서는 \textbf{hit와 miss에서 필요한 memory access 횟수}를 먼저 세는 것이 핵심이다.
\end{solutionframe}

% 11 Physics - 문제
\begin{problemframe}{11}{Physics}
한 변의 길이가 $a$, 저항이 $R$인 정사각형 도선고리가 속력 $v$로 오른쪽으로 움직이며, 종이면 안쪽 방향의 균일한 자기장 $B$가 존재하는 영역으로 들어가고 있다. 고리가 자기장 영역에 \textbf{부분적으로 들어가는 동안}, self-inductance를 무시할 때 운동을 방해하는 자기력의 크기를 구하시오.
\end{problemframe}

% 11 Physics - 정답
\begin{solutionframe}{11}{Physics}
\answerbox{$\displaystyle F=\frac{B^2a^2v}{R}$}
\solutionlabel
자기선속 증가율은 $\dd\Phi/\dd t=Bav$이므로
\[
\mathcal E=Bav,
\qquad
I=\frac{Bav}{R},
\qquad
F=IaB=\frac{B^2a^2v}{R}.
\]
Lenz's law에 따라 이 힘은 운동 반대쪽으로 작용한다.
\end{solutionframe}

% 12 Chemistry - 문제
\begin{problemframe}{12}{Chemistry}
Octahedral complex $\ce{[Fe(CN)6]^4-}$와 $\ce{[FeF6]^4-}$에서 Fe의 oxidation state는 모두 $+2$이다. $\ce{CN^-}$는 strong-field ligand, $\ce{F^-}$는 weak-field ligand라 할 때, 두 complex의 unpaired electron 수를 순서대로 쓰시오.
\end{problemframe}

% 12 Chemistry - 정답
\begin{solutionframe}{12}{Chemistry}
\answerbox{$\displaystyle 0,\ 4$}
\solutionlabel
$\ce{Fe^{2+}}$는 $d^6$이다. Strong-field $\ce{CN^-}$에서는 low-spin $t_{2g}^6$가 되어 unpaired electron이 $0$개이다. Weak-field $\ce{F^-}$에서는 high-spin $t_{2g}^4e_g^2$가 되어 unpaired electron이 $4$개이다. 즉 전자는 같아도 \textbf{ligand-field splitting의 크기}가 spin state를 바꾼다.
\end{solutionframe}

% 13 Biology - 문제
\begin{problemframe}{13}{Biology}
\normalsize
산소와 기질이 충분한 isolated mitochondria에 protonophore인 2,4-dinitrophenol (DNP)을 첨가하였다. Proton motive force를 $\Delta p$라 할 때, 첨가 직후의 변화를 가장 잘 나타낸 것을 고르시오.
\begin{simchoices}
\item $\mathrm{O_2}$ consumption $\downarrow$, ATP synthesis $\uparrow$, $\Delta p\uparrow$
\item $\mathrm{O_2}$ consumption $\uparrow$, ATP synthesis $\downarrow$, $\Delta p\downarrow$
\item $\mathrm{O_2}$ consumption $\downarrow$, ATP synthesis $\downarrow$, $\Delta p\uparrow$
\item $\mathrm{O_2}$ consumption $\rightarrow$, ATP synthesis $\uparrow$, $\Delta p\downarrow$
\end{simchoices}
\end{problemframe}

% 13 Biology - 정답
\begin{solutionframe}{13}{Biology}
\answerbox{(2)}
\solutionlabel
DNP는 inner mitochondrial membrane을 가로질러 proton을 운반해 proton motive force를 소실시키는 \textbf{uncoupler}이다. ATP synthase가 이용할 gradient가 줄어 ATP synthesis는 감소한다. 반면 electron transport chain에 걸리던 ``back pressure''가 줄어 respiration과 $\mathrm{O_2}$ consumption은 증가한다.
\end{solutionframe}

% 14 Mathematics - 문제
\begin{problemframe}{14}{Mathematics}
Congruence
\[
x^2\equiv1\pmod{105}
\]
를 만족하는 residue class $x\pmod{105}$의 개수를 구하시오.
\end{problemframe}

% 14 Mathematics - 정답
\begin{solutionframe}{14}{Mathematics}
\answerbox{$\displaystyle 8$}
\solutionlabel
$105=3\cdot5\cdot7$이고 서로 다른 odd primes의 곱이다. 각 prime $p\in\{3,5,7\}$에 대해
\[
x^2\equiv1\pmod p
\]
의 해는 $x\equiv\pm1$ 두 개이다. Chinese Remainder Theorem에 의해 세 선택을 독립적으로 조합할 수 있으므로 해의 개수는
\[
2^3=8.
\]
\end{solutionframe}

% 15 Computer Engineering - 문제
\begin{problemframe}{15}{Computer Engineering}
어떤 block code의 minimum Hamming distance가 $d_{\min}=7$이다. 이 code가 항상 검출할 수 있는 오류 bit 수의 최댓값과 항상 정정할 수 있는 오류 bit 수의 최댓값을 순서대로 쓰시오.
\end{problemframe}

% 15 Computer Engineering - 정답
\begin{solutionframe}{15}{Computer Engineering}
\answerbox{검출 $6$ bits, 정정 $3$ bits}
\solutionlabel
Minimum distance가 $d_{\min}$이면 최대 $d_{\min}-1$개의 오류를 항상 detect할 수 있고,
\[
\left\lfloor\frac{d_{\min}-1}{2}\right\rfloor
\]
개의 오류를 항상 correct할 수 있다. $d_{\min}=7$이므로 각각 $6$과 $3$이다.
\end{solutionframe}

% 16 Physics - 문제
\begin{problemframe}{16}{Physics}
두 별 A와 B를 모두 blackbody로 근사하며 반지름은 서로 같다. B의 blackbody spectrum에서 peak wavelength가 A의 정확히 절반이라면, B의 total luminosity는 A의 몇 배인가? Wien's law와 Stefan-Boltzmann law를 사용할 수 있다.
\end{problemframe}

% 16 Physics - 정답
\begin{solutionframe}{16}{Physics}
\answerbox{$\displaystyle \frac{L_B}{L_A}=16$}
\solutionlabel
Wien's law $\lambda_{\max}T=\text{const.}$에 의해 peak wavelength가 절반이면
\[
T_B=2T_A.
\]
반지름이 같으므로 Stefan-Boltzmann law $L=4\pi R^2\sigma T^4$에서
\[
\frac{L_B}{L_A}=\left(\frac{T_B}{T_A}\right)^4=2^4=16.
\]
\end{solutionframe}

% 17 Chemistry - 문제
\begin{problemframe}{17}{Chemistry}
2025 Nobel Prize in Chemistry는 metal-organic frameworks (MOFs)의 개발에 수여되었다. 전형적인 MOF의 crystalline porous framework를 구성하는 두 종류의 building block을 쓰시오.
\end{problemframe}

% 17 Chemistry - 정답
\begin{solutionframe}{17}{Chemistry}
\answerbox{Metal ions/clusters + multidentate organic linkers}
\solutionlabel
MOF는 metal ion 또는 metal cluster를 node로 하고, 여러 배위점을 갖는 organic ligand/linker가 이들을 연결하여 만든 porous coordination network이다. Building block을 바꾸면 pore size와 chemical environment를 조절할 수 있어 gas storage, $\ce{CO2}$ capture, catalysis 등에 활용된다.
\sourceinfo{Nobel Prize in Chemistry 2025}{https://www.nobelprize.org/prizes/chemistry/2025/press-release/}
\end{solutionframe}

% 18 Biology - 문제
\begin{problemframe}{18}{Biology}
두 유전자 $A$와 $B$의 recombination frequency가 $20\%$이다. Heterozygote의 haplotype이 $AB/ab$이고 이를 $aabb$와 testcross하였다. 자손이 $aabb$일 확률을 구하시오.
\end{problemframe}

% 18 Biology - 정답
\begin{solutionframe}{18}{Biology}
\answerbox{$\displaystyle 40\%$}
\solutionlabel
Recombinant gametes $Ab$, $aB$의 총 비율이 $20\%$이므로 각각 $10\%$이다. 따라서 nonrecombinant gametes $AB$, $ab$는 총 $80\%$, 각각 $40\%$이다. Tester는 항상 $ab$ gamete를 주므로 $aabb$ 자손은 heterozygote가 $ab$를 줄 때 생기며 그 확률은 $40\%$이다.
\end{solutionframe}

% 19 Mathematics - 문제
\begin{problemframe}{19}{Mathematics}
2026 Abel Prize를 받은 수학자로, Mordell conjecture를 증명하여 ``genus가 $1$보다 큰 $\Q$ 위의 smooth projective curve는 rational point를 유한개만 갖는다''는 결과를 정리로 만든 인물의 이름을 쓰시오.
\end{problemframe}

% 19 Mathematics - 정답
\begin{solutionframe}{19}{Mathematics}
\answerbox{Gerd Faltings}
\solutionlabel
Gerd Faltings는 1983년 Mordell conjecture를 증명했고, 오늘날 이 결과를 Faltings' theorem이라 부른다. 그는 arithmetic geometry에 강력한 도구를 도입하고 Mordell 및 Lang의 오래된 Diophantine conjectures를 해결한 업적으로 2026 Abel Prize를 받았다.
\sourceinfo{The Abel Prize 2026 - Gerd Faltings}{https://abelprize.no/article/2026/gerd-faltings-awarded-2026-abel-prize}
\end{solutionframe}

% 20 Computer Engineering - 문제
\begin{problemframe}{20}{Computer Engineering}
Charles H. Bennett와 Gilles Brassard는 quantum information science의 토대를 세운 업적으로 2025 ACM A.M. Turing Award를 받았다. 두 사람이 1984년에 제안한 대표적인 quantum key distribution protocol의 이름을 쓰시오.
\end{problemframe}

% 20 Computer Engineering - 정답
\begin{solutionframe}{20}{Computer Engineering}
\answerbox{BB84 protocol}
\solutionlabel
BB84는 Bennett와 Brassard가 1984년에 제안한 quantum key distribution (QKD) protocol이다. 서로 호환되지 않는 quantum basis를 사용하며, 도청자가 unknown quantum state를 측정하면 상태를 교란하여 오류율 상승으로 도청 흔적을 검출할 수 있다는 것이 핵심이다.
\sourceinfo{2025 ACM A.M. Turing Award - Bennett and Brassard}{https://www.nsf.gov/cise/updates/pioneers-quantum-information-science-recognized-2025-acm-am}
\end{solutionframe}

\end{document}
